\section{工程与监控}

\subsection{RLHF Pipeline Timeline}

把 RLHF pipeline 拆成阶段，便于团队协作：

\begin{itemize}
    \item 数据构建周：instruction + comparison data 采集
    \item 奖励建模周：重复训练 + calibration + validation
    \item 策略优化周：PPO 或 DPO 迭代 + KL/entropy tuning
    \item 质量检查周：benchmark + red team + human eval
    \item 监控与 rollout：gradual deployment + drift monitoring
\end{itemize}

\begin{importantbox}{Action Timeline}
每个阶段都定义清晰的 deliverable：数据构建结果、奖励模型 blueprints、PPO checkpoint、审查报告、部署决策。这样团队才知道什么时候可以推进下一步。
\end{importantbox}

\subsection{Operational Monitoring}

部署后必须持续监控 hallucination、rejection、latency。

\begin{figure}[H]
\centering
\includegraphics[width=0.8\textwidth]{slide-01.jpg}
\caption{Slide 1 把 RLHF pipeline 与监控指标同时展示，强调 latency 与 rejection 的追踪。}
\end{figure}

\begin{knowledgebox}{观察盘的指标设计}
推荐同时展示：Reward score drift、human rejection rate、token latency、deploy coverage。不同指标组合可以帮助团队快速判断是 model drift 还是 data drift。
\end{knowledgebox}

\begin{table}[H]
\centering
\begin{tabular}{p{0.3\textwidth} p{0.6\textwidth}}
\toprule
指标 & 说明 \\
\midrule
Reward score drift & 用 rolling window 比较 current vs baseline reward distribution \\
Reject ratio & 审查反馈的拒绝率，用于评估 guardrail 灵敏度 \\
Latency P99 & 保障 transformer + verifier pipeline 不超出 SLO \\
Human override count & 人工干预次数，用来 calibrate automation trust \\
\bottomrule
\end{tabular}
\caption{部署后观察盘核心指标}
\end{table}

\subsection{Governance Checklist}

\begin{table}[H]
\centering
\begin{tabular}{p{0.2\textwidth} p{0.35\textwidth} p{0.28\textwidth}}
\toprule
控制点 & 核心问答 & 输出 \\
\midrule
Reward model approval & 是否经过 blind eval validation? & Approval checklist \\
Policy guardrails & 是否有 automatic rejection for forbidden topics? & Guardrail rules \\
Deployment gating & 是否逐步扩容、扣留异常? & Rollout policy \\
\bottomrule
\end{tabular}
\caption{RLHF 的治理检查表}
\end{table}

\subsection{Drift Detection 架构}

一旦模型上线，需要检测两类 drift：prompt drift（用户输入变化）与 distribution drift（输出行为变化）。

\begin{itemize}
    \item 在 logging pipeline 中捕获 prompt signature（hash + metadata）
    \item 用 rolling window 统计 reward score distribution
    \item 提供 alert，当 KL divergence vs zscore 超阈值时触发 human review
    \item 建立 `drift ticket`，包含 snapshot prompt, reward score, output diff
\end{itemize}

\begin{knowledgebox}{Drift detection 不只是 anomaly detection}
它需要 context-aware thresholds：不同任务（summaries vs math）对 reward score 的 variance 本身就不同。缺乏 context-specific threshold，会造成大量假告警。
\end{knowledgebox}

\subsection{本章小结}

工程层面要把 timeline、监控和治理结合起来，监控 dashboard 和 checklist 是部署后信任的重要基础。
\subsection{Ops Playbook}

\begin{enumerate}
    \item \textbf{Observe}：聚合 reward scores、rejection logs、prompt hash。
    \item \textbf{Assess}：用 evidence matrix 与 counterfactual prompt 检查 drift。
    \item \textbf{Act}：让 AI 推荐 low-risk 缓解动作并生成 rollback plan，人工批准后执行。
    \item \textbf{Document}：把决策写回 checklist / runbook，供下一轮 evaluation 使用。
\end{enumerate}

\begin{table}[H]
\centering
\begin{tabular}{p{0.18\textwidth} p{0.4\textwidth} p{0.3\textwidth}}
\toprule
阶段 & 输出 & 核心参与者 \\
\midrule
Observe & Reward + guardrail logs & SRE + Data Eng \\
Assess & Evidence scorecard + drift report & Alignment + QA \\
Act & Rollback plan/mitigation action & Product + Legal \\
Document & Runbook entry + retrospective & Ops + PM \\
\bottomrule
\end{tabular}
\caption{Ops Playbook 的可执行交付物}
\end{table}

\begin{knowledgebox}{Ops runbook 的价值}
把每次 Act 往返写成 runbook 条目，不仅可以 guarantee 复现，也在 board review 时提供透明的 audit trail。
\end{knowledgebox}

\begin{warningbox}{AI 主要做 triage，决策仍需人类把关}
即使 AI 建议某条规则，仍要在高风险线路上让人类审核，防止自动化误判导致不可逆后果。
\end{warningbox}

\subsection{本章小结}

工程层面要把 timeline、监控、治理、drift detection 与 ops playbook 串成闭环，才能在灰度/全量阶段维持 alignment 信任。

